
%\begin{center}
%\hrule \medskip % Upper rule
%\large{Incinération et environnement}
%\normalsize 
%\vspace{0.3cm}

%Unité 3\\ 

%\medskip\hrule 
%\bigskip
%\end{center}

\chapter{Incinération et environnement}

Cette unité de formation vise à fournir des informations sur :
\begin{itemize}
\item l'environnement ;
\item les impacts de l'incinération des déchets dangereux (DD).
\end{itemize}

Ce module aborde les grands thèmes de l’environnement (air, eau, sols, cadre de vie et nuisances, déchets) selon deux axes, celui des généralités sur les problèmes de chaque milieu et celui des aspects spécifiques des usines d’incinération de déchets dangereux, appelées dans notre métier UIDD. 
\vspace{0.3cm}

D’autre part sont abordées de façon succincte la réglementation environnementale et la norme ISO 14001.


\section{Impacts atmosphériques}

\subsection{Air et pollution atmosphérique} 

\begin{figure}[H]
    \centering
    \includegraphics[width=\textwidth]{Images/M3/emissions_pollutants.png}
\end{figure}

L’air est un mélange de plusieurs gaz et il est principalement composé d’azote (78\%), d’oxygène (21\%) et d’argon (0,9\%). Les 0,1\% restants sont constitués de multiples composés en très faible concentration.
Certains polluants, présents en quantités infimes, peuvent entraîner des effets nuisibles à la santé. Plus la durée d’exposition au polluant est importante, plus ces effets risquent d’être aggravés. On parle de pollution atmosphérique lorsque certains composants sont présents dans l’air à des concentrations supérieures à la normale et qu’ils peuvent entraîner une modification du milieu naturel ou occasionner une gêne pour les êtres vivants, notamment l’être humain. Les phénomènes de pollution atmosphérique peuvent être d’échelle locale, régionale ou mondiale.
\\

En milieu urbain, la circulation automobile et les activités industrielles sont les principaux responsables d’une pollution locale limitée à quelques dizaines de kilomètres autour des agglomérations.
Certains phénomènes, comme les pluies acides ou la pollution à l’ozone, dépassent le cadre local et concernent une ou plusieurs régions (50 à 100 km, voire plus).
\\
D’autres phénomènes plus complexes affectent l’atmosphère terrestre dans sa globalité. Les effets de cette pollution globale sont ressentis sur le long terme (élévation de la température de l’air, trou dans la couche d’ozone, etc.).
Le schéma de la page précédente montre les principaux phénomènes de pollution atmosphérique.


\subsection{Qualité de l'air}

L’Organisation mondiale de la santé (OMS) et l’Union européenne (UE) ont défini des 
valeurs-seuil (concentration maximale par unité de temps) à respecter pour de nombreux polluants atmosphériques.
En France, plus de 550 sites, généralement situés dans les grandes agglomérations, mesurent en continu la concentration de plusieurs polluants. Le tableau suivant donne le résultat de la surveillance de quelques polluants en 1995 :


\begin{figure}[H]
    \centering
    \includegraphics[width=\textwidth]{Images/M3/air_quality.png}
\end{figure}

Remarque : Les termes généraux sont définis plus en détail dans les fiches de définition jointes à la fin de chaque unité. Triez ces feuilles à la fin du dossier par ordre alphabétique.



\subsection{Dispersion des polluants atmosphériques}


Les rejets gazeux industriels sont généralement canalisés et extraits par des cheminées qui ont pour rôle de disperser les polluants et d’éviter leur retombée directe au niveau du sol. 
\\
\\
\\
La bonne dispersion de ces polluants dépendra :

\begin{itemize}

   \item des conditions d’émission : vitesse et température des gaz au niveau du rejet, hauteur de la cheminée ;
   \item des conditions météorologiques.
   \end{itemize}
   
   
La réglementation relative aux UIDD impose de calculer la hauteur de cheminée et la vitesse d’extraction en prenant en compte :
\begin{itemize}
   \item les concentrations admissibles au sol en poussières, en acide chlorhydrique et en dioxyde de soufre;
   \item la présence éventuelle d’autres sources de rejets de polluants atmosphériques à proximité;
   \item la présence d’obstacles proches (bâtiments) qui pourraient gêner la dispersion des fumées.
Les schémas ci-dessous montrent deux modes de dispersion des rejets d’une cheminée.
 
\end{itemize}

\begin{figure}[H]
    \centering
    \includegraphics[width=\textwidth]{Images/M3/normal_diffusion.png}
\end{figure}

Diffusion normale
(Source : La Qualité de l'air, notre patrimoine vital, (La qualité de l'air, notre patrimoine vital) Ministère français de l'Environnement, 1989)

La bonne dispersion des polluants peut être perturbée ou la concentration en polluants peut être augmentée en fonction de conditions météorologiques particulières :

\begin{itemize}
    
\item lorsque la température de l’air, à quelques centaines de mètre d’altitude, est supérieure à la température du sol;
\item lorsque l’air chaud des grandes agglomérations entraîne vers le centre ville l’air pollué des zones industrielles périphériques plus fraîches.
\end{itemize}

\begin{figure}[H]
    \centering
    \includegraphics[width=\textwidth]{Images/M3/inv_temp.png}
\end{figure}

Inversion de température
(Source : La Qualité de l'air, notre patrimoine vital, (La qualité de l'air, notre patrimoine vital) Ministère de l'Environnement, 1989)

\subsection{Rejets atmosphériques de l'UIDD} 

Rappel : composés polluants issus de l'incinération avant traitement des fumées : 
\begin{itemize}
\item les poussières sont produites en grande quantité par la combustion des déchets ménagers;
\item l’acide chlorhydrique (HCl) provient principalement de la combustion de solvants chlorés et d’emballages plastiques;
\item on trouve également, mais dans des proportions faibles par rapport à leur potentiel polluant, les composés suivants : métaux lourds, composés organiques volatils (COV), oxydes d’azote (NOX), dioxyde de soufre (SO2), dioxines et furanes;
\item le gaz carbonique (CO2). Rappelons que le gaz carbonique n’est pas un polluant, car il est présent en grande quantité dans l’atmosphère. L’incinération génère du gaz carbonique qui, à quantité égale (c’est-à-dire incinération d’une tonne de déchets ou stockage d’une tonne de déchets), contribue bien moins à l’augmentation de l’effet serre que le méthane, les CFC ou les oxydes d’azote.

\end{itemize} 

Le traitement des fumées des UIDD, qui répond aux normes de rejet de l’arrêté du 25 janv. 1991 établies par le ministère de l’Environnement, permet une épuration très importante des différents polluants.
Le tableau suivant représente les quantités en poids des différents polluants émis dans les gaz de combustion, valeurs exprimées en kilogramme (kg) pour une tonne de déchets incinérés, valeurs émises avant et après traitement :

\textbf{Exemples de poids de polluants émis dans les fumées des gaz. 
(hypothèse calculée sur la base de 6000 m3 par tonne) :}

\begin{figure}[H]
    \centering
    \includegraphics[width=\textwidth]{Images/M3/pollutant_emitted.png}
\end{figure}

\textbf{Dioxines : les rejets étant très faibles, il est préférable d’exprimer les valeurs de rejets gazeux pour 100000 t de déchets incinérés, soit :}
\begin{itemize}
\item en sortie d'incinérateur : 0,6 à 6 g pour 100 000 t de déchets incinérés ;
\item après traitement : 0,06 g pour 100 000 t de déchets ou 60 milligrammes pour 100 000 t de déchets.
   
\end{itemize}


Une autre manière de comparer les valeurs des rejets gazeux est de présenter les quantités en poids des différents polluants émis dans les gaz de combustion en milligrammes (mg) par mètre cube de gaz rejeté à l’atmosphère.


\begin{figure}[H]
    \centering
    \includegraphics[width=\textwidth]{Images/M3/pollutant_emitted2.png}
\end{figure}

Ainsi une nouvelle usine traitant 100000 t/an de déchets dangereux rejetterait dans 
l’atmosphère moins de 1/20 g de dioxine par an.
\vspace{0.3cm}

Rappelons brièvement les procédés de traitement les plus utilisés pour piéger ces polluants que nous étudierons dans les prochains modules (modules 9 et 10) :


\begin{figure}[H]
    \centering
    \includegraphics[width=\textwidth]{Images/M3/pollutant_emitted3.png}
\end{figure}

\subsection{Réglementation des rejets atmosphériques} 

La réglementation applicable aux UIDD impose des valeurs limites pour un grand nombre de polluants contenus dans les rejets. Les principaux polluants (HCl, SOx, poussières, monoxyde de carbone (CO), Carbone Organique Total (COT) et Oxydes d’Azote(NOx)) doivent être mesurés en continu.
\vspace{0.3cm}

Les mesures et le contrôle en continu des valeurs des polluants permettent à l’exploitant de vérifier le bon fonctionnement de son installation de traitement des fumées, de réajuster en permanence la quantité de réactifs nécessaires à la neutralisation et à l’épuration des gaz pollués (voir le détail de la réglementation dans les fiches de définition). Ainsi est limitée la quantité de polluants rejetés à l’atmosphère.
D’autres composés, comme les métaux lourds, ne sont mesurés qu’une fois par an.



\begin{figure}[H]
    \centering
    \includegraphics[width=\textwidth]{Images/M3/pollutant_mesured.png}
\end{figure}


Le tableau suivant présente quelques valeurs issues de la réglementation en vigueur : 

\begin{figure}[H]
    \centering
    \includegraphics[width=\textwidth]{Images/M3/regulation.png}
\end{figure}

Les mesures de dioxines sont délicates à réaliser. C’est la raison pour laquelle notre groupe a décidé d’équiper le Creed (Centre de recherche pour l’environnement, l’énergie et le déchet, situé à Limay, dans les Yvelines) d’un laboratoire de mesures de dioxines. Ce laboratoire réalise en France des prélèvements d’échantillon de fumées et des analyses.

\begin{figure}[H]
    \centering
    \includegraphics[width=0.5\textwidth]{Images/M3/dioxine_mesure.png}
\end{figure}

D’autre part, en France, il existe une taxe parafiscale payable par les exploitants des usines sur leurs émissions polluantes à l’atmosphère. Cette taxe concerne les émissions suivantes :
\begin{itemize}
 \item oxydes de soufre (SO2, etc.) ;
 \item oxydes d'azote (NOX, etc.) ;
 \item acide chlorhydrique (HCl);
 \item poussière.
\end{itemize}

L’Ademe (Agence de l’environnement et de la maîtrise de l’énergie) utilise les fonds collectés de cette taxe pour financer des actions et des investissements liés à la prévention et à la réduction de la pollution atmosphérique en France.


\subsection{Contribution des rejets polluants des UIDD à la pollution de l'air}

Les sources de pollution de l’air sont multiples (la circulation urbaine, le chauffage individuel, les activités industrielles…). Des bilans ont été établis et montrent comment les polluants émis par les UIDD, contribuent à la pollution de l’air.
\vspace{0.3cm}

Le tableau ci-dessous résume l’importance relative à ces rejets.


\begin{figure}[H]
    \centering
    \includegraphics[width=\textwidth]{Images/M3/pollutant_emitted4.png}
\end{figure}


\section{Impacts sur l'eau}

\subsection{Généralités sur l'eau} 

Omniprésente sur terre, on trouve l’eau sous cinq formes naturelles différentes : 
\begin{itemize}
   
\item les eaux de surface (rivières, lacs, océans);
\item les eaux souterraines (infiltration, nappe phréatique);
\item  les dépôts de glace et de neige (forme solide);
\item l’atmosphère (vapeur, forme gazeuse);
\item  la biosphère (toutes les formes de vie sur terre).

\end{itemize}


L’eau circule entre ces différentes formes suivant le cycle naturel de l'eau.
Nécessaire à la vie, l’eau se retrouve à tous les échelons de la chaîne alimentaire et constitue le principal constituant des êtres vivants (70\% du poids du corps humain).

\newpage


\subsection{Caractérisation des eaux naturelles}

\textbf{Eaux souterraines}
\vspace{0.3cm}

Les nappes phréatiques peuvent être alimentées directement par l’infiltration des eaux 
de ruissellement ou de façon indirecte lorsqu’elles sont séparées du sol par une couche imperméable. Suivant la nature géologique du sol, les eaux des nappes peuvent avoir des 
compositions chimiques très variables. Ces eaux sont généralement caractérisées par l’absence d’oxygène dissous, une grande pureté bactériologique, une faible turbidité (non transparence de l’eau), une température et une composition chimique constante. L’eau des nappes est souvent potable. En cas de pollution accidentelle, toutes ces caractéristiques peuvent être fortement modifiées et il devient alors très difficile de récupérer la pureté originelle de la nappe.


\vspace{0.3cm}
\textbf{Eau de surface}
\vspace{0.3cm}

Elles ont pour origine des nappes profondes (sources, ruisseaux) ou des eaux de ruissellement issues des bassins versants.
La composition chimique des eaux dépend de la nature des terrains traversés par l’eau 
durant son parcours. Par échange à la surface eau-atmosphère, ces eaux se chargent en gaz dissous (oxygène, azote, gaz carbonique). 
Parmi les caractéristiques des eaux de surfaces, il faut noter :



\begin{itemize}

\item la présence d’oxygène dissous;
\item une concentration non négligeable de matières en suspension aux tailles très variables;
\item des matières organiques d’origine naturelle provenant de la décomposition des organismes végétaux et animaux;
\item la présence de plancton, d’algues, de poissons et autres.

    
\end{itemize}

\vspace{0.3cm}
\textbf{Sources de pollution de l'eau}
\vspace{0.3cm}

Il y a pollution de l’eau s’il y a modification de ses caractéristiques et si ce phénomène risque de remettre en cause les usages qui en sont faits (potabilité, survie des poissons, etc.). Cette pollution est d’origine naturelle ou humaine. La pollution naturelle peut être due à des précipitations importantes ou aux conséquences d’une longue période de sécheresse. Les principales sources de pollution d’origine humaine peuvent être classées comme suit :
\begin{itemize}
  
\item les eaux résiduaires urbaines, qui doivent être traitées en station d’épuration collective avant rejet ;
\item les effluents industriels, qui sont généralement traités sur le site industriel ou bien sur la station d’épuration collective lorsque la pollution générée est comparable à celle engendrée par les ménages ;
\item la pollution d’origine agricole (par exemple, les nitrates présents dans les nappes phréatiques suite à l’utilisation excessive d’engrais).
  
\end{itemize}

\subsection{Les principaux aspects de la pollution de l'eau}

\vspace{0.3cm}
\textbf{Contenu en oxygène}
\vspace{0.3cm}

Dans les eaux de surface, l’oxygène joue un rôle fondamental dans le maintien de la vie aquatique et dans l’auto-épuration. L’oxygène représente environ 35\% des gaz dissous dans l’eau. De multiples réactions chimiques peuvent entraîner la consommation d’oxygène dissous. De plus, lorsque la température augmente, la solubilité de l’oxygène diminue et les bactéries, plus nombreuses, consomment plus d’oxygène. En dessous de 5 mg/l, les poissons sont en danger.

\vspace{0.3cm}
\textbf{Pollution organique}
\vspace{0.3cm}

D’origine naturelle, urbaine ou industrielle, la pollution organique est caractérisée par la 
présence de matières organiques biodégradables (feuilles mortes…) ou synthétiques (solvants, hydrocarbures, détergents…). Ces matières organiques ont tendance à se dégrader 
naturellement en consommant de l’oxygène.
Deux paramètres fréquents permettent de caractériser l’altération en oxygène par pollution organique : la DCO et la DBO (voir fiches de définition).

\vspace{0.3cm}
\textbf{Matières en suspension (MES)}
\vspace{0.3cm}

Les petits éléments non dissous dans l’eau diminuent la transparence de l’eau (on dit aussi que la turbidité de l’eau augmente). Les MES comprennent des minéraux, des matières 
organiques de taille importante, des virus, des bactéries, etc. En diminuant le rayonnement solaire à travers l’eau, les MES perturbent et diminuent l’activité des organismes qui constituent la base de la chaîne alimentaire. En se déposant au fond des eaux calmes, ces solides peuvent emprisonner des micro-polluants. Enfin, ils peuvent également provoquer l’asphyxie des poissons par colmatage des ouïes.

\vspace{0.3cm}
\textbf{Neutralité}
\vspace{0.3cm}

Le pH est un paramètre chimique mesuré sur une échelle de 0 à 14. Il permet d’évaluer le degré d’acidité d’une solution.

\begin{figure}[H]
    \centering
    \includegraphics[width=\textwidth]{Images/M3/neutrality.png}
\end{figure}

Lorsque le pH est inférieur à 6 ou supérieur à 9, la vie des poissons est mise en danger, voire impossible. Si le pH d’une solution est égal à 7, on dit que la solution est neutre.

\vspace{0.3cm}
\textbf{Temperature}
\vspace{0.3cm}

C’est un facteur important de la vie aquatique. Une élévation de la température due à des 
activités industrielles peut avoir de graves conséquences sur le métabolisme des organismes vivants. Au-dessus de 22°C, on observe la mort de certaines espèces aquatiques, comme les truites.


\vspace{0.3cm}
\textbf{Métaux lourds}
\vspace{0.3cm}

Même à faible concentration, leur danger provient surtout de leur capacité à s’accumuler dans les chaînes alimentaires.

Notons également deux autres types de pollution non détaillés ici : la pollution azotée et la pollution phosphorée, significatives pour de nombreuses industries et pour l’agriculture, mais ne concernant pas directement l’activité des UIDD.


\vspace{0.3cm}
\textbf{Qualité des eaux de surface}
\vspace{0.3cm}

En France, plus de 1500 points de mesure permettent de suivre la qualité des cours d’eau. Selon les résultats des mesures effectuées sur de nombreux paramètres, à chaque cours d’eau est attribué un indice de qualité. 

Il existe cinq indices de qualité :

\begin{itemize}
   
\item Classe 1A : eaux exemptes de pollution aptes à satisfaire les usages les plus exigeants (eau potable, baignade, etc.);
\item  Classe 1B : d’une qualité légèrement moindre, ces eaux peuvent néanmoins satisfaire tous les usages;
\item  Classe 2 : qualité passable : nécessite un traitement poussé pour être potable. Le poisson y vit normalement mais sa reproduction est aléatoire;
\item  Classe 3 : juste apte à l’irrigation et à l’industrie (refroidissement, par exemple). La subsistance de la vie piscicole y est aléatoire;
\item  Hors classe : considérée comme inapte à la plupart des usages et constituant une menace pour la santé publique et l’environnement.

\end{itemize}



Le tableau ci-dessous donne des exemples de teneurs limites pour chaque classe de qualité :

\begin{figure}[H]
    \centering
    \includegraphics[width=\textwidth]{Images/M3/class_water.png}
\end{figure}

Le ministère de l’Environnement fixe un objectif de qualité pour chaque cours d’eau. Par exemple, une rivière actuellement classée 2 (passable) peut avoir un objectif de qualité de classe 1B. Ces objectifs de qualité servent de référence aux administrations pour fixer les valeurs limites de rejet de chaque installation. 


\subsection{Impact sur les nappes d'eaux souterraines}

La pollution d’une nappe souterraine est souvent due à une contamination du sol situé à proximité. Lorsqu’un site industriel se trouve à proximité d’une nappe sensible, un réseau de surveillance de la qualité de l’eau souterraine est exigé par l’administration. Des piézomètres, sortes de mini puits forés dans le sol, permettent d’effectuer des prélèvements de l’eau de la nappe et de vérifier par analyse la qualité des eaux souterraines
En France, la qualité des eaux potables est réglementée et très suivie. L’Organisation mondiale de la santé (OMS), la Communauté économique européenne (CEE) et la France ont défini des valeurs recommandées, ainsi que des valeurs maximales à ne pas dépasser pour près de 100 paramètres chimiques. Le tableau ci-dessous reprend les valeurs limites recommandées pour divers métaux lourds :


\begin{figure}[H]
    \centering
    \includegraphics[width=\textwidth]{Images/M3/drinking_water.png}
\end{figure}

\subsection{Rejets liquides de UIDD}
\vspace{0.3cm}

\textbf{Les eaux de process}
Les équipements ou activités qui peuvent être à la source de rejets liquides sont :
\begin{itemize}
\item le traitement des fumées, principalement les procédés « type humide ». Les purges des laveurs font l’objet d’un traitement d’eau spécifique qui permet de corriger le pH, de piéger les métaux lourds, de décanter les matières en suspension et de filtrer les boues. Les rejets après traitement peuvent contenir des sels de chlorure en quantité plus ou moins importante;
\item  la production de vapeur (dans les usines avec valorisation énergétique) : les éluats de la chaîne de production d’eau déminéralisée, les purges de la chaudière;
\item  le refroidissement des mâchefers : l’eau utilisée pour refroidir les mâchefers doit être collectée et traitée avant rejet (sur les installations récentes, cette eau est recyclée);
\item  d’autres circuits auxiliaires (eau de refroidissement de pompes alimentaires, etc.);
\item  les rejets des unités de traitement de désodorisation.
\end{itemize}

Après traitement éventuel (bac de décantation, correction de pH), tous ces effluents sont 
rejetés vers le milieu naturel (rivière, fleuve) ou vers un réseau d’assainissement public raccordé à une station d’épuration collective. Certaines usines ne produisent pas d’effluents liquides.


Lorsque les effluents de l’UIDD sont rejetés vers le milieu naturel, les valeurs limites fixées par la Drire (Direction régionale de l’industrie, de la recherche et de l’environnement), organisme d’État chargé de vérifier l’application des normes de rejets, sont adaptées à la nature du cours d’eau et à son indice de qualité. 
Dans tous les cas, l’UIDD doit respecter les prescriptions de l’arrêté ministériel du 20 septembre 2002, notamment les paramètres suivants :


\begin{figure}[H]
    \centering
    \includegraphics[width=\textwidth]{Images/M3/discharges.png}
\end{figure}
\vspace{0.3cm}

\textbf{Lorsque les effluents sont rejetés vers un réseau d’assainissement public}, 
les valeurs de rejet sont généralement fixées en accord avec l’exploitant du réseau ou de la station municipale de traitement dans une convention de raccordement. Les prescriptions peuvent être moins contraignantes pour les MES et la DCO que celles indiquées ci-dessus. 
Exemple : quelques prescriptions de rejets de l’UIOM de Saran


\begin{figure}[H]
    \centering
    \includegraphics[width=\textwidth]{Images/M3/discharges2.png}
\end{figure}

Toutefois, pour les polluants qui ne peuvent être éliminés par la station municipale, comme les métaux lourds, les prescriptions de rejet de la convention sont identiques à celle de rejet vers le milieu naturel.
Dans les deux cas, un ou plusieurs traitements d’eau au préalable sont nécessaires pour respecter les prescriptions de rejet (neutralisation, décantation, etc.).


\textbf{Les eaux pluviales}
Les eaux de pluie sont généralement collectées séparément des eaux de process. Elles font l’objet d’un traitement minimal (débourbeur, déshuileur) avant rejet vers le milieu naturel. Parfois, un bassin de rétention permet de retenir toutes les eaux de pluie afin d’éviter les pollutions accidentelles du milieu récepteur (par exemple, en cas de déversement accidentel de fuel, acide, chaux). 
Enfin, certains points de collecte des eaux de pluie peuvent être raccordés au réseau des eaux de process pour traitement (exemple : eaux de pluie à proximité du poste de réception et stockage de la chaux).

\textbf{Fréquence des mesures}
Fréquence des mesures
Selon les arrêtés préfectoraux, l’exploitant doit effectuer des mesures quotidiennes, hebdomadaires, mensuelles ou tous les six mois. Parfois, il est exigé un enregistrement en continu des débits, du pH et de la température des rejets.

\section{Impacts sur les sols}

\subsection{Les fonctions du sol}

Comme l’air et l’eau, le sol est une ressource naturelle essentielle. S’il est pollué, le sol peut porter atteinte à la vie humaine, en affectant les récoltes ou, plus insidieusement, en contaminant les végétaux, les aliments ou l’eau. 
Des trois milieux (air, eau, sol), le sol est probablement le moins bien connu. 
 
En tant que milieu naturel, le sol assure quatre fonctions différentes :

\begin{itemize}

\item un rôle de tampon et de régulation hydrologique, limitant les crues et la sécheresse des cours d’eau;
\item  un rôle de filtre : rétention de composés toxiques préservant ainsi la nappe phréatique;
\item  un rôle de transformation microbienne en dégradant les substances toxiques organiques;
\item  un rôle protecteur pour les nombreux organismes qui y vivent, nécessaires aux écosystèmes.

   
\end{itemize}

L’infiltration des polluants lors de la contamination des eaux souterraines est explicitée à la figure ci-dessous.

\begin{figure}[H]
    \centering
    \includegraphics[width=\textwidth]{Images/M3/underground_water.png}
\end{figure}

L’approvisionnement en eau industrielle d’un site de production se fait par de multiples moyens (pompage en rivière, approvisionnement urbain, par exemple). Parfois, une part souvent importante provient de prélèvements directs dans la nappe phréatique à l’aplomb du site.
C’est dire l’importance qu’il convient d’attacher au sol et au sous-sol d’un site industriel : toute pollution a une répercussion directe sur les approvisionnements du site et des communes, proches ou lointaines.

\vspace{0.3cm}
La réglementation française prévoit qu’autour de chaque captage d’eau potable, doit être définie une zone de protection dans laquelle des mesures spéciales sont prises afin d’éviter tout risque de pollution de la source.


\subsection{La pollution des sols}

La pollution des sols résulte de l’activité humaine. Elle se caractérise par la dégradation des capacités régulatrices du sol (protection des nappes souterraines, alimentation des végétaux, etc.) et est susceptible de provoquer une nuisance ou un risque pour les personnes ou pour l’environnement. 
\vspace{0.3cm}

Les conséquences de cette pollution peuvent être importantes, par exemple lorsqu’une ressource en eau potable devient inutilisable ou lorsque l’usage du sol pour la culture devient limité ou impossible.

\begin{figure}[H]
    \centering
    \includegraphics[width=\textwidth]{Images/M3/underground_water2.png}
\end{figure}


\vspace{0.3cm}
\textbf{La pollution du sol peut s’effectuer selon deux formes différentes :}
\begin{itemize}

\item pollution diffuse : ce type de pollution peut être lié à la pollution atmosphérique; les retombées atmosphériques, d’origine urbaine, industrielle ou routière, s’infiltrent dans les sols (par exemple, les métaux lourds ou les retombées des pluies acides). Les pratiques agricoles peuvent également conduire à la pollution du sol et des nappes phréatiques (ce peut être le cas des engrais, des produits pour la protection des plantes ou des boues de station d’épuration qui sont lessivés par les pluies et entraînés dans les nappes phréatiques). Les stockages de déchets, à même le sol, non conformes à la réglementation peuvent également être une source de pollution diffuse (lixivation de déchets placés en décharge lorsque le terrain n’est pas suffisamment imperméable).
\item pollution ponctuelle : généralement localisées à un emplacement délimité, les causes de pollution ponctuelles peuvent être nombreuses. Citons, par exemple, les déversements accidentels ou volontaires de produits polluants, les accidents de la circulation (camions citernes contenant des produits dangereux), les dépôts sauvages de déchets, les fuites de produits chimiques.
\end{itemize}


\subsection{Le sol et ses polluants} 

\textbf{Parmi les polluants du sol et sous-sol on peut distinguer principalement :}

\begin{itemize}
\item les métaux lourds : au niveau du sol, ils peuvent être piégés en se transformant au contact d’autres éléments chimiques. Mais ils peuvent également traverser le sol et se retrouver dans les nappes souterraines. Ils peuvent enfin se retrouver dans les plantes ou légumes en agricole. Les concentrations en métaux lourds dans les sols sont généralement beaucoup plus élevées à proximité des zones urbaines ou industrielles. Toutefois, les sols contiennent naturellement des métaux lourds à des concentrations très variables, dans des rapports pouvant varier de 1 à 5, et exceptionnellement de 1 à 50.
\item les micro-polluants organiques : les hydrocarbures, solvants, PCB (polycycliques aromatiques chlorés) et les produits phytosanitaires à base de chlore; ces produits peuvent contaminer la chaîne alimentaire.
\item les éléments pathogènes : bactéries, virus, parasites et autres micro-organismes qui peuvent contaminer les sols après déjections animales, épandage de boues ou amendements organiques non fermentés.
\end{itemize}

En France, il n’existe pas de valeurs-seuil à respecter pour la contamination des sols par les différents types de polluants. Une réglementation est en cours de préparation.




\subsection{La prévention des impacts sur les sols : Le cas des UIDD } 


Les rejets atmosphériques des UIDD engendrent de faibles retombées particulaires de 
métaux lourds, dioxines et autres polluants contenus dans les fumées d’incinération. Ces 
retombées sont négligeables car diffuses. Des analyses régulières des sols et végétaux 
permettent de vérifier l’absence de concentration et d’accumulation dans les sols à proximité de l’UIDD.
Les aspects liés aux risques de pollution des sols par les résidus de l’incinération (résidus d’épuration à des fumées, REFIDD, et mâchefers) seront abordés dans les prochains modules.


\vspace{0.3cm} 


Les UIDD doivent éliminer les risques de déversement accidentel d’hydrocarbures, de réactifs chimiques ou d’autres produits dangereux pour le sol et l’eau. Pour éviter ces accidents, des moyens de rétention, des fosses étanches (cuvettes de rétention en béton) doivent être prévus et installés aux endroits de stockage et de manipulation (fosse de stockage des mâchefers, fosse de réception des déchets dangereux, rétention des stockage de déchets liquides …).


\section{Cadre de vie et nuisances}


La qualité du cadre de vie prend une part de plus en plus importante dans nos sociétés modernes. Respirer un air sain, boire une eau potable et de qualité, s’alimenter avec des aliments sans risques pour la santé ne sont pas les seules préoccupations des ménages. D’autres aspects comme les odeurs, le bruit, la qualité paysagère de l’environnement sont à prendre en compte lors de la conception des UIDD et dans leur mode d’exploitation.


\subsection{Les Odeurs et les Composés Organiques Volatiles (COV)}

Le monde qui nous entoure est plein d'odeurs. Une multitude de molécules différentes mélangées à l’air que nous respirons, généralement à des concentrations extrêmement faibles, sont à l’origine des odeurs que nous percevons. Pour que ces molécules soient odorantes, il faut qu’elles réagissent avec la muqueuse nasale et que notre cerveau l’interprète comme une sensation olfactive. Cette dernière dépend de la nature de la substance odorante, et de sa concentration dans l’air que nous respirons.


\vspace{0.3cm}

\textbf{Odeurs \& COV générés par l'UIDD}
\vspace{0.3cm}

\textbf{Plusieurs sources de COV peuvent être détectées :}


\begin{itemize}

\item	les fosses de réception des déchets et les stockages de déchets liquides : sur les UIDD pour éviter les odeurs et protéger les travailleurs, les COV sont captés et brûlés avec l’air primaire qui sert à la combustion des déchets ;
\item		les zones de pré-traitement des déchets (broyage, pompage…) : ces opérations sont généralement réalisées dans des bâtiments et l’air de ces zones est traité avant rejet à l’atmosphère.
\item	d’autres sources moins importantes comme les odeurs des mâchefers.

\end{itemize}

Souvent, une bonne maîtrise des procédés et une organisation orientée vers la prévention permettent de réduire les nuisances potentielles. 
La réglementation relative aux odeurs impose généralement que l’UIDD ne doit pas 
générer des odeurs qui incommodent le voisinage. Parfois, lorsqu’un système de désodorisation est installé, l’arrêté d’exploitation précise des seuils de concentration à respecter en limite de propriété pour certains composés odorants. 


\subsection{Bruit}

Plus de 50\% des Français considèrent le bruit comme l'une des principales sources de nuisance.

\vspace{0.3cm}
\textbf{Qu'est-ce que le bruit ?}
\vspace{0.3cm} 

Le bruit est un phénomène acoustique produisant une sensation considérée comme désagréable ou gênante. Tout son est dû à une variation de la pression régnant dans l’atmosphère, engendrée par une source sonore. Cette variation de pression, même infime, est appelée pression acoustique. Le niveau sonore s’exprime en décibels (dB).
Lorsqu’un bruit double d’intensité, le niveau sonore exprimé en décibels augmente de 3 dB. Si ce bruit triple d’intensité, le niveau sonore augmente de 5 dB.
Le tableau ci-dessous donne des exemples de niveaux de bruit.

\begin{figure}[H]
    \centering
    \includegraphics[width=\textwidth]{Images/M3/noise.png}
\end{figure}

\textbf{Bruit généré par l'UIDD}
\vspace{0.3cm}

\textbf{Les sources possibles de bruit dans une UIDD sont multiples. On peut noter les exemples suivants :}

\begin{itemize}
   
 \item camions, machines, alarmes ; 
 \item turbocompresseurs, ventilateurs, systèmes à vapeur, tours de refroidissement d'air, soupapes de décharge ;
\item broyeurs, cribles, machines vibrantes.

\end{itemize}


 
Le bruit du poste de travail dans l'UIDD ne doit pas dépasser 80 dB. Certains locaux et équipements plus bruyants doivent être insonorisés et les opérateurs doivent utiliser des dispositifs de protection spécifiques.
Sur le plan environnemental, le bruit à l'extérieur du site est réglementé et ne doit pas dépasser des seuils variables selon les sites. Fréquemment, selon la réglementation, le niveau sonore généré par l'usine en bordure de propriété ne doit pas dépasser le bruit de fond extérieur de plus de 5 dB en journée en semaine (ce qui correspond à x3 le niveau sonore), et de 3 dB la nuit , les week-ends et jours fériés (x2 le niveau sonore). Cette notion de bruit supplémentaire généré par une installation est appelée émergence.

\vspace{0.3cm}

Par ailleurs, le niveau sonore maximum admissible en bordure de propriété est défini en fonction du type de zone dans laquelle se situe le site UIDD. Le tableau suivant donne les exigences applicables pour les installations situées dans une zone dominée par des opérations commerciales et industrielles.


\begin{figure}[H]
    \centering
    \includegraphics[width=\textwidth]{Images/M3/noise2.png}
\end{figure}


\subsection{Intégration paysagère}  

Les installations industrielles modifient le paysage dans lequel elles s’insèrent. Visibles de près ou de loin par les personnes extérieures, il est important qu’elles soient maintenues dans un état optimal de propreté compte tenu du service qu’elles rendent. 

Parmi les idées contribuant à une bonne conception d’une nouvelle usine, retenons les principes suivants : 

\begin{itemize}
\item optimiser les volumes et les surfaces occupés par l’UIDD en les réduisant tout en permettant une exploitation optimale;
\item retenir une architecture de prestige qui soit à l’image des bonnes performances environnementales de l’UIDD;
\item choisir des matériaux et des coloris qui facilitent l’intégration de l’usine dans le paysage.

\end{itemize}
 
La contribution de l’exploitant à la réduction des nuisances visuelles peut être de maintenir un bon état de propreté extérieure (nettoyage régulier adapté, maintien de la « fraîcheur » de la peinture d’origine, balayage des aires extérieures…).

\begin{figure}[H]
    \centering
    \includegraphics[width=0.7\textwidth]{Images/M3/image_limay.png}
\end{figure}


\section{Impacts liés aux déchets}

\subsection{L’incinération : un traitement permettant de réduire les impacts potentiels}


L’incinération est, à l’heure actuelle, le moyen de traitement des déchets dangereux le mieux adapté pour respecter l’environnement.
 
L’incinération permet de réduire l’impact potentiel des déchets, en les réduisant considérablement :

 \begin{itemize}
\item d’une proportion de 1 à 10 en volume environ;
\item  d’une proportion de 1 à 3 en poids environ.

 \end{itemize}

\subsection{Matériaux et déchets issus du traitement} 

\vspace{0.3cm}
\textbf{Résidus d'épuration des fumées (REFIDD)}
\vspace{0.3cm}

La combustion d’une tonne de déchets dangereux produit environ 50 à 100 kg de résidus d’épuration des fumées. Les REFIDD récupérés après dépoussiérage ou lavage des gaz sont des polluants potentiels du milieu naturel. Ils peuvent transmettre en cas de mauvaise gestion de leur élimination leur pollution au sol et à l’eau et atteindre ainsi les espèces végétales animales et bien sûr les êtres humains.

La composition des REFIDD est variable suivant la technique de neutralisation des gaz utilisée. Tous ces résidus présentent une forte solubilité. D’autre part on constate des teneurs en métaux lourds élevées notamment pour le plomb.
L’eau est le vecteur principal d’une pollution accidentelle. Une lixivation accidentelle solubilise les polluants de ces résidus puis les dissémine dans le sol. Pour les rendre plus inertes il faut leur faire subir un traitement qui empêche que les polluants puissent être emportés par l’eau. Ainsi avant stockage en Centre de Stockage de Déchets Ultimes (CSDU) les REFIDD sont stabilisés en les solidifiant dans des réactifs solides et liquides de façon à obtenir un matériau solide comparable à du béton. Ces obligations sont décrites dans différents textes réglementaires.


\vspace{0.3cm}
\textbf{Les mâchefers}
\vspace{0.3cm}

La combustion d’une tonne de déchets dangereux produit environ 200 à 300 kg de résidus solides à la sortie des fours. Ces mâchefers se caractérisent par :
\begin{itemize}

\item	des matériaux inertes non solubles ;
\item		des sels (solubles) ;
\item		de l’eau (18 à 30\% suivant les installations) ;
\item		des métaux ferreux (environ 10\%) qui peuvent être captés en sortie de four par tri magnétique.
\end{itemize}

Le taux d’imbrûlés doit être inférieur à 5\%. Le pH des mâchefers est généralement voisin de 11 ou 12.

Les mâchefers déferraillés sont considérés comme des résidus ultimes et doivent donc être stockés en CSDU, tel quel ou après stabilisation.


Remarque :
Concernant les mâchefers d’UIOM (Usine d’Incinération d’Ordures Ménagères), la circulaire ministérielle du 9 mai 1994 encourage et réglemente la valorisation des mâchefers. Si l’on souhaite valoriser les mâchefers, il faut procéder à une caractérisation pour les classer dans l’une des trois catégories possibles :
 \begin{itemize}
\item comme valorisable (intitulée V) : ils peuvent être utilisés sous certaines conditions, par exemple en tant que remblais ou de comblement comme en sous-couche routière;
\item comme mâturable (intitulée M) : une étape de maturation ultérieure est nécessaire pour que le mâchefer puisse éventuellement être valorisé;
\item  comme stockable (intitulée S) : dans ce cas, seule la mise en décharge de classe 2 est envisageable.
  \end{itemize}

Cette caractérisation résulte de plusieurs analyses effectuées durant six mois. Après cette caractérisation, une analyse mensuelle doit être réalisée. Chaque mois on effectue la moyenne des sept dernières analyses ce qui permet de déterminer la catégorie des mâchefers (V, M ou S). 

Concernant les mâchefers d’UIDD, cette possibilité de valorisation n’est pas offerte. Ils doivent obligatoirement aller en CSDU.



\section{La réglementation et l'environnement}

\subsection{Introduction}

La réglementation française sur la protection de l’environnement s’articule autour de 2 dispositifs complémentaires :
\begin{itemize}
\item des lois visant chaque catégorie de nuisance ou de milieu naturel : c’est ainsi que l’on distingue principalement les lois sur l’eau, l’air, les déchets, le bruit;
\item des lois recherchant des solutions globales pour des activités ou installations à caractère polluant : c’est le cas de la réglementation sur les installations classées.

\end{itemize}


\newpage
\textbf{La législation des Installations classées :}
\vspace{0.3cm}

La loi n° 76-663 du 19 juillet 1976 définit le cadre réglementaire applicable aux installations qui sont soumises soit à déclaration, soit à autorisation, et fournit la nomenclature officielle qui s’applique à chaque type d’activité, et un arrêté préfectoral d’autorisation d’exploitation.
Ces installations sont suivies régulièrement par un inspecteur de la DRIRE. Celui-ci peut, s’il l’estime nécessaire, compléter le cadre législatif applicable à l’installation par des dispositions spécifiques, sous la forme de prescriptions complémentaires.

\subsection{La réglementation applicable aux UIDD}

Plusieurs dizaines de textes réglementaires concernent complètement ou partiellement les activités d’incinération de déchets ménagers. L’organigramme ci-dessous présente de façon simplifiée les principaux textes applicables aux UIDD.

\begin{figure}[H]
    \centering
    \includegraphics[width=\textwidth]{Images/M3/applicable_regulation.png}
\end{figure}

\section{ISO 14001 : un modèle de gestion environnementale}

\subsection{La gestion des contraintes environnementales}

Les activités de traitement de déchets ménagers peuvent engendrer des impacts non 
négligeables sur l’environnement. Les chapitres précédents ont abordé dans les grandes lignes ces impacts sur l’air, l’eau, le sol ainsi que les nuisances potentielles sur le cadre de vie. 
Pour réduire ces impacts, une réglementation lourde et complexe prévoit de nombreuses contraintes. Respecter cette réglementation n’est pas chose facile : 

\begin{itemize}
\item les équipements de traitement de déchets et de dépollution doivent être adaptés et conformes. 
\item l’exploitant d’une UIDD doit organiser toutes les actions nécessaires : optimisation des procédés, surveillance et mesures régulières sur site, choix de filières de valorisation ou d’élimination conformes à la réglementation, transmission de rapports aux administrations et aux clients, etc.
\item  la réglementation évolue sans cesse, en parallèle avec la prise de conscience de 
problèmes jusqu’alors mal connus (exemple : pollution de l’air, dioxines, pollution des sols) et dans le cadre d’une réglementation internationale et européenne de plus en plus contraignante.

\end{itemize}


De plus, les attentes de la population en matière de santé publique sont croissantes : 
qualité de l’air, intolérance aux odeurs et au bruit, craintes sur la contamination de la chaîne alimentaire. Souvent, ces attentes dépassent le cadre de la réglementation et il convient d’anticiper sur les évolutions futures des textes de loi. Pour rassurer le voisinage et les associations, de nombreux exploitants doivent engager des actions de communication et de concertation sur la qualité environnementale des UIDD. 

Les enjeux sont également financiers et juridiques : la responsabilité des entreprises est grande en cas de pollution (obligation d’indemniser et de dépolluer) et des UIDD peuvent être fermées en cas de risque d’atteinte à la santé.
 
En résumé, pour gérer au mieux les contraintes environnementales d’une UIDD il faut :

\begin{itemize}
\item respecter scrupuleusement la réglementation;
\item  maîtriser au mieux ses impacts et les risques de pollution;
\item communiquer avec tous les acteurs concernés dans un souci de transparence;
\item anticiper sur les évolutions rapides.
 \end{itemize}

\subsection{ISO 14000 : un outil de management standardisé}

En 1996, l’Organisation internationale de normalisation (ISO, Insurance services office) a édité un modèle de gestion des problèmes d’environnement qui puisse être reconnu par tous mondialement : les normes ISO 14000.
Parmi ces normes, ISO 14001 est un référentiel d’organisation à respecter pour les entreprises qui souhaitent gérer au mieux les contraintes environnementales décrites ci-dessus.

Ces normes ne sont pas obligatoires mais correspondent à une démarche volontaire. En appliquant les règles d’ISO 14001, l’entreprise prend un triple engagement :

\begin{itemize}
\item respecter la réglementation environnementale;
\item  chercher à protéger l’environnement et prévenir les risques de pollution;
\item  améliorer en continu la gestion des problèmes d’environnement.
\end{itemize} 

L’intérêt d’ISO 14001 est d’être une norme internationalement reconnue : une fois la norme respectée par l’entreprise, celle-ci peut demander à un organisme extérieur de venir vérifier à travers un audit que son organisation est conforme. L’entreprise est alors certifiée ISO 14001 et peut l’annoncer à ses clients en France et à l’international.
Ce certificat ISO 14001 est une preuve de compétence en matière de gestion environnementale. C’est un atout sérieux pour gagner de nouveaux marchés d’exploitation.


\subsection{le contenu ISO 14001}

\textbf{Pour mettre en place ISO 14001, il faut :}
\begin{itemize}
\item   effectuer un examen initial détaillé de sa situation environnementale : textes réglementaires applicables, inventaires des impacts environnementaux, sensibilité du milieu extérieur, maîtrise des activités;
\item  choisir des objectifs d’amélioration à travers un programme chiffré. Pour une UIDD, cela peut par exemple consister à réduire ses rejets atmosphériques ou recycler les eaux usées ou bien optimiser la valorisation des produits (mâchefers, compost);
\item  rédiger une politique environnementale adaptée aux objectifs d’amélioration;
\item  formaliser sous forme de procédures et documents opératoires les pratiques à respecter (exemple : réception des déchets, conduite du traitement de fumées, analyses des rejets en eau, etc.);
\item   former le personnel aux impacts environnementaux;
\item   assurer une veille réglementaire afin de détecter les évolutions des textes liés à l’environnement;
\item   effectuer des audits internes réguliers pour vérifier que l’organisation est conforme aux 
dispositions écrites;
\item   effectuer des revues de Direction pour améliorer l’organisation en continu.

\end{itemize}

La notion d’amélioration continue est présentée dans le schéma ci-dessous :

\begin{figure}[H]
    \centering
    \includegraphics[width=0.5\textwidth]{Images/M3/continuous_improvment.png}
\end{figure}

\section{Index}

\subsection{Incinération de déchets dangereux en France}




\begin{figure}[H]
    \centering
    \includegraphics[width=\textwidth]{Images/M3/SARPI_france.png}
\end{figure}

\subsection{Réglementation applicable à l'UIDD : Arrêté du 20 septembre 2002}

Réglementation applicable aux UIDD : principaux points de l’Arrêté Ministériel Incinération du 20 septembre 2002 (entrée en vigueur le 28 décembre 2005).

Conditions de mesure
Toutes les concentrations sont exprimées en mg/Nm3 sec ramenés à 11\% d’O2
Les moyennes sur une demi-heure et les moyennes sur dix minutes sont déterminées pendant la période de fonctionnement effectif ( à l’exception des phases de démarrage et d’extinction, lorsqu’aucun déchet n’est incinéré) à partir des valeurs mesurées après soustraction de l’intervalle de confiance à 95 p. 100 sur chacune de ces mesures. Cet intervalle de confiance ne doit pas dépasser les pourcentages suivants des valeurs limites d’émission :

\begin{figure}[H]
    \centering
    \includegraphics[width=0.5\textwidth]{Images/M3/regulation_UIDD.png}
\end{figure}

Les moyennes journalières sont calculées à partir de ces moyennes validées.

Pour qu’une moyenne journalière soit valide, il faut que, pour une même journée, pas plus de cinq moyennes sur une demi-heure n’aient dû être écartées pour cause de mauvais fonctionnement ou d’entretien du système de mesure en continu. Dix moyennes journalières par an peuvent être écartées au maximum pour cause de mauvais fonctionnement ou d’entretien du système de mesure en continu.

Condition de combustion
Les installations d’incinération sont conçues, équipées, construites et exploitées de manière à ce que, même dans les conditions les plus défavorables que l’on puisse prévoir, les gaz résultant du processus soient portés, après la dernière injection d’air de combustion, d’une façon contrôlée et homogène, à une température de 850 ° C pendant deux secondes, mesurée à proximité de la paroi interne ou en un autre point représentatif de la chambre de combustion défini par l’arrêté préfectoral d’autorisation. Le temps de séjour devra être vérifié lors des essais de mise en service. S’il s’agit de déchets dangereux ayant une teneur en substances organiques halogénées, exprimée en chlore, supérieure à 1 p. 100, la température doit être amenée à 1100 ° C pendant au moins deux secondes.
La température doit être mesurée en continu.

Chaque ligne d’incinération est équipée d’au moins un brûleur d’appoint, lequel doit s’enclencher automatiquement lorsque la température des gaz de combustion tombe en dessous de 850 ° C ou de 1100 ° C, selon le cas, après la dernière injection d’air de combustion. Ces brûleurs sont aussi utilisés dans les phases de démarrage et d’extinction afin d’assurer en permanence la température de 850 °C ou de 1100 ° C, selon le cas, pendant lesdites phases et aussi longtemps que des déchets non brûlés se trouvent dans la chambre de combustion.

\newpage
\textbf{Normes d'émission}
\vspace{0.3cm}

\begin{figure}[H]
    \centering
    \includegraphics[width=\textwidth]{Images/M3/emissions_standards_pollutants.png}
\end{figure}


\textbf{Respect des Valeurs Limite d’Emissions}
\vspace{0.3cm}

Les valeurs limites d’émission sont respectées si :

\begin{itemize}
\item aucune des moyennes journalières mesurées ne dépasse les limites d’émission fixées à l’article 17 pour le monoxyde de carbone et pour les poussières totales, les substances organiques à l’état de gaz ou de vapeur exprimées en carbone organique total (C.O.T.), le chlorure d’hydrogène, le fluorure d’hydrogène, le dioxyde de soufre et les oxydes d’azote ;
\item  aucune des moyennes sur une demi-heure mesurées pour les poussières totales, les substances organiques à l’état de gaz ou de vapeur exprimées en carbone organique total, le chlorure d’hydrogène, le fluorure d’hydrogène, le dioxyde de soufre et les oxydes d’azote ne dépasse les valeurs limites définies à l’article 17 ;
\item  aucune des moyennes mesurées sur la période d’échantillonnage prévue pour le cadmium et ses composés ainsi que le thallium et ses composés, le mercure et ses composés, le total des autres métaux (Sb+As+Pb+Cr+Co+Cu+Mn+Ni+V), les dioxines et furannes, ne dépasse les valeurs limites définies à l’article 17.
\item  95 p. 100 de toutes les moyennes mesurées sur dix minutes pour le monoxyde de carbone sont inférieures à 150 mg/m3 ; ou aucune mesure correspondant à des valeurs moyennes calculées sur une demi-heure au cours d’une période de 24 heures ne dépasse 100 mg/m3.
\end{itemize}

\textbf{Demande Chimique en Oxygène (DCO) et Demande Biologique en Oxygène (DBO)}

\begin{itemize}
 \item Demande chimique en oxygène (DCO) : c’est la quantité d’oxygène consommée dans un litre d’eau par la totalité les matières organiques oxydables. Cette consommation est proportionnelle à la quantité de matière organique chimiquement oxydable, et donc à son niveau de pollution.
 \item Demande biologique en oxygène (DBO) : c’est la quantité d’oxygène nécessaire à la destruction ou à la dégradation des matières organiques d’une eau par les micro-organismes du milieu. Pour des raisons pratiques, on utilise beaucoup la DBO5 (mesure de la DBO en 5 jours). Cette consommation est proportionnelle à la quantité de matière organique biologiquement oxydable, et donc à son niveau de pollution.
\end{itemize}

\newpage
\textbf{Métaux lourds}
\vspace{0.3cm}

Métaux lourds : les déchets dangereux contiennent de nombreux métaux qui pour la plupart se retrouvent dans les résidus solides de combustion (mâchefers cendres volantes). Des métaux sont également présents dans les fumées majoritairement sous forme de poussières et aussi sous forme gazeuse.
L’antimoine (Sb), l’arsenic (As), le plomb (Pb), le chrome (Cr), le cobalt (Co), le cuivre (Cu), le manganèse (Mn), le nickel (Ni), le vanadium (V), ainsi que le mercure (Hg), le cadmium (Cd) et le thallium (Tl).

\vspace{0.3cm}

Rappelons que les métaux lourds les plus dangereux pour la santé humaine sont le cadmium, le plomb et le mercure.

\vspace{0.3cm}

\textbf{Pourquoi ces métaux sont-ils considérés comme "lourds" ?} 
\vspace{0.3cm} 

Les éléments traces sont des éléments chimiques présents en très faibles quantités dans les différents milieux (sol eau air). 
Parmi ces éléments traces certains sont nécessaires aux organismes vivants (plantes animaux l’homme) dans de très faibles quantités : il s’agit des oligo-éléments (exemple : zinc, cuivre, manganèse, bore, molybdène, cobalt, sélénium). Remarquons qu’à de trop grandes quantités et au-delà de certains seuils ces oligo-éléments peuvent devenir toxiques.
D’autres éléments traces (plomb, cadmium, mercure, nickel, arsenic…) ne sont pas utiles aux organismes et au-delà d’un certain seuil peuvent devenir toxiques pour le milieu naturel et l’homme. Ils sont désignés communément par le terme de « métaux lourds ». Le terme « métaux lourds » généralement utilisé en langage environnemental est souvent un abus de langage car tous ces éléments ne sont pas chimiquement des métalloïdes (exemple : l’arsenic). 
En chimie on distingue les métaux dits lourds (à cause de leur poids effectif), des métaux légers (calcium, aluminium…), des métaux nobles (or, platine…) ou encore des métaux rares.


\vspace{0.3cm}

\textbf{Dioxines et furanes} 
\vspace{0.3cm}

\begin{figure}[H]
    \centering
    \includegraphics[width=0.4\textwidth]{Images/M3/PCDD.png}
\end{figure}

Dioxines et furanes : ces deux familles de polluants se forment principalement lors de la combustion de matières contenant du chlore comme les plastiques ou certains solvants (on trouve naturellement des composés chlorés dans le bois et le charbon). 
Le terme « dioxine » regroupe 75 types de dioxines et 135 types de furanes. Ces dioxines ont pour faculté de s’accumuler dans les graisses et de s’éliminer très difficilement.

\vspace{0.3cm}
\newpage
\textbf{Formule chimique :}
\vspace{0.3cm}

\textbf{Polychlorodibenzodioxine (PCDD) et polychlorodibenzofuranes (PCDF)}

La voie alimentaire constitue à 95\% la seule source d’exposition. L’homme apparaît bien moins sensible à leur exposition que les animaux. Une seule catégorie (la 2 3 7 8 TCDD) est classée substance cancérogène pour l’homme mais avec un risque faible. Dans l’attente de mieux connaître ces risques les recommandations et normes font valoir le principe de sécurité.
\vspace{0.3cm}

Les opérations de prélèvements et d’analyses sont longues délicates et coûteuses (environ 5000 € par opération). L’évaluation normalisée de la charge toxique porte sur 17 isomères (les plus toxiques) et correspond à une moyenne pondérée des 17 toxicités d’importance inégale.
\vspace{0.3cm}

En France la contribution des Usines d’Incinération d’Ordures Ménagères (UIOM) aux émissions industrielles totales de dioxines a été évaluée à 43\% d’après l’Ademe en 1996. Les UIOM Veolia contribuent à hauteur d’un tiers aux émissions des UIOM françaises. Néanmoins, depuis 1996, de nombreuses UIOM anciennes, hors normes, ont été fermées, et la plupart des usines en fonctionnement respecte désormais la norme. Les UIDD sont soumis à cette norme depuis juillet 2000 et avait anticipé son entrée en vigueur. Leur contribution aux émissions de dioxines est donc très faible.
\vspace{0.3cm}

La nouvelle directive Européenne sur l’incinération des déchets (Wastes Incinération Directive WID), transposée en droit français par l’Arrêté Ministériel du 20 septembre 2002, impose à tous les incinérateurs et co-incinérateur de déchets une norme de rejet en dioxines et furanes de 0,1 I-TEQOTAN ng/Nm3 sec à 11\%O2.
Cet arrêté est applicable aux usines nouvelles immédiatement et à toutes les usines à partir du 28 décembre 2005.
\vspace{0.3cm}

Remarque : 0.1 ng/m3 = 10-10 g/m3 = 0.1 milliardième de gramme par m3. Une nouvelle usine d’incinération qui traiterait 100 000 t/an de déchets dangereux produirait moins de 1/20 de gramme de dioxines toxiques durant l’année.
 


\vspace{0.3cm}

\textbf{Symboles chimiques}
\vspace{0.3cm}

Toute matière vivante (minérale, végétale, organique, qu'elle soit liquide, solide ou gazeuse) est formée d'éléments de base identifiés par un symbole internationalement reconnu.

\vspace{0.3cm}

Exemple :	N : azote		O : oxygène

\vspace{0.3cm}

Chacun de ces éléments de base se caractérise par une particule élémentaire qui lui est propre : son atome. Ce dernier se compose de plus petites unités, les électrons et les protons, qui détermineront les conditions d’association entre les éléments de base pour former d’autres substances.
\vspace{0.3cm}

Ainsi, dans certaines conditions de température et de pression, il y a compatibilité entre l’azote et l’oxygène pour former une nouvelle substance (NO2) ou oxyde d’azote. Comme il s’agit d’un regroupement d’éléments de base, la particule élémentaire du NO2 s’appelle molécule.
\vspace{0.3cm}

La plupart des substances, produits ou matériaux que nous utilisons chaque jour sont des regroupements d’éléments de base, mais il y a certaines exceptions : l’or, le cuivre, l’aluminium peuvent avoir des degrés de pureté assez élevés et sont des éléments de base.